% ============================================================
%  Chapter 2 -- Perfect Sets and Connected Sets, 2.43-2.47
% ============================================================
\rsection{Perfect Sets}

\begin{signpost}[Two results, one of them a warning]
Theorem 2.43 is a second, self-contained proof that $\R$ is uncountable, using
compactness instead of Cantor's diagonal --- and unlike 2.14 it needs no
discussion of binary expansions.

\smallskip
Section 2.44 is the real business. The Cantor set is the standing
counterexample of the subject: compact, perfect, uncountable, yet containing no
interval and (as Chapter 11 will show) of measure zero. Every intuition of the
form ``a set that big must contain an interval'' dies here. Expect to meet it
again in Exercises 2.17--2.18, in Exercise 6.6 (integration against the Cantor
set), and in Chapter 11.
\end{signpost}

\ritem{2.43}{Theorem}
\emph{Let $P$ be a nonempty perfect set in $\R^k$. Then $P$ is uncountable.}

\begin{proof}
Since $P$ has limit points, $P$ must be infinite. Suppose $P$ is countable, and
denote the points of $P$ by $\mathbf{x}_1,\mathbf{x}_2,\mathbf{x}_3,\dots$. We
shall construct a countable collection $\{V_n\}$ of neighborhoods, as
follows.\kn{The strategy is Cantor's diagonal in disguise: assume a list of all
of $P$, then trap a point that the list misses. Where 2.14 built the missing
object digit by digit, here it is caught by a nested sequence of compact sets
--- and 2.36 guarantees something survives.}

Let $V_1$ be any neighborhood of $\mathbf{x}_1$. If $V_1$ consists of all
$\mathbf{y} \in \R^k$ such that $|\mathbf{y}-\mathbf{x}_1| < r$, then the
closure $\overline{V}_1$ of $V_1$ is the set of all $\mathbf{y} \in \R^k$ such
that $|\mathbf{y}-\mathbf{x}_1| \le r$.

Suppose $V_n$ has been constructed, so that $V_n \cap P$ is not empty. Since
every point of $P$ is a limit point of $P$, there is a neighborhood $V_{n+1}$
such that (i) $\overline{V}_{n+1} \subset V_n$, (ii)
$\mathbf{x}_n \notin \overline{V}_{n+1}$, (iii) $V_{n+1} \cap P$ is not
empty.\kn{Clause (ii) is the diagonal step --- it evicts the $n$th listed point
at stage $n$. Clause (iii) is where \emph{perfect} is used: because
$\mathbf{x}$ is a limit point of $P$, there are always more points of $P$
nearby to keep the construction alive.} By (iii), $V_{n+1}$ satisfies our
induction hypothesis, and the construction can proceed.

Put $K_n = \overline{V}_n \cap P$. Since $\overline{V}_n$ is closed and
bounded, $\overline{V}_n$ is compact. Since $\mathbf{x}_n \notin K_{n+1}$, no
point of $P$ lies in $\bigcap_{n=1}^\infty K_n$. Since $K_n \subset P$, this
implies that $\bigcap_{n=1}^\infty K_n$ is empty. But each $K_n$ is nonempty,
by (iii), and $K_n \supset K_{n+1}$, by (i); this contradicts the Corollary to
Theorem 2.36.
\end{proof}

\begin{center}
\begin{tikzpicture}[x=1mm, y=1mm]
  \draw[thin] (0,0) circle (20mm);
  \draw[thin] (5,3) circle (11mm);
  \draw[thin] (8,5) circle (5.5mm);
  % evicted points
  \node[ptfill, minimum size=2.6pt] at (-13,-7) {};
  \node[lbl, anchor=east] at (-23,-7) {$\mathbf{x}_1$: inside $V_1$,};
  \node[lbl, anchor=east] at (-23,-10.5) {evicted by $V_2$};
  \draw[thin] (-22,-7.5) -- (-15,-7);
  \node[ptfill, minimum size=2.6pt] at (-1,-2) {};
  \draw[thin] (-1,-4) -- (6,-22.5);
  \node[lbl, anchor=north] at (8,-23.5) {$\mathbf{x}_2$: inside $V_2$, evicted by $V_3$};
  % survivor
  \node[ptopen] at (9,6) {};
  \draw[thin] (11,7.5) -- (26,13);
  \node[lbl, anchor=west] at (27,15) {something of $P$ survives};
  \node[lbl, anchor=west] at (27,11.5) {every stage (2.36) ---};
  \node[lbl, anchor=west] at (27,8) {yet the list was supposed};
  \node[lbl, anchor=west] at (27,4.5) {to contain all of $P$};
  % circle labels
  \node[mlbl, anchor=south east] at (-11,13) {$V_1$};
  \node[mlbl, anchor=south east] at (0,9) {$V_2$};
  \node[mlbl] at (8,1.7) {$V_3$};
\end{tikzpicture}
\end{center}
\figcap{Stage $n{+}1$ shrinks inside stage $n$ and evicts the $n$th listed
point; perfectness keeps each stage nonempty. Compactness forces a survivor,
which the list therefore missed --- Cantor's diagonal, drawn with balls
instead of digits.}

\begin{center}
\begin{tikzpicture}[x=1mm, y=1mm]
  \node[rmbox, text width=27mm] (c0) at (0,22)
    {suppose $P$ countable:\\list as $\mathbf{x}_1,\mathbf{x}_2,\dots$};
  \node[rmbox, text width=29mm] (c1) at (44,22)
    {build $V_{n+1} \subset V_n$:\\evicts $\mathbf{x}_n$, meets $P$};
  \node[rmbox, text width=27mm] (c2) at (87,22)
    {$K_n = \overline{V}_n \cap P$:\\compact, nested, $\ne\emptyset$};
  \draw[rmarrow] (c0) -- (c1);
  \draw[rmarrow] (c1) -- node[lbl, anchor=south, inner sep=2.5pt]
    {\;} (c2);
  \node[rmbox, text width=27mm] (c3) at (87,0)
    {2.36: $\bigcap K_n \ne \emptyset$};
  \draw[rmarrow] (c2) -- (c3);
  \node[rmbox, text width=29mm] (c4) at (44,0)
    {but every $\mathbf{x}_n$ evicted:\\$\bigcap K_n \cap P = \emptyset$};
  \draw[rmarrow] (c2) -- (c4);
  \node[rmbox, text width=26mm] (c5) at (0,0) {contradiction};
  \draw[rmarrow] (c4) -- (c5);
  \draw[rmarrow] (c3) -- (c4);
\end{tikzpicture}
\end{center}
\figcap{The flow: the list drives the construction, perfectness keeps every
stage alive, and compactness (2.36) collides with total eviction.}

\ritem{}{Corollary}
\emph{Every interval $[a,b]$ $(a<b)$ is uncountable. In particular, the set of
all real numbers is uncountable.}

\ritem{2.44}{The Cantor Set}
The set which we are now going to construct shows that there exist perfect sets
in $\R^1$ which contain no segment.

Let $E_0$ be the interval $[0,1]$. Remove the segment
$\bigl(\tfrac13,\tfrac23\bigr)$, and let $E_1$ be the union of the intervals
\[ \Bigl[0,\tfrac13\Bigr], \qquad \Bigl[\tfrac23,1\Bigr]. \]
Remove the middle thirds of these intervals, and let $E_2$ be the union of the
intervals
\[ \Bigl[0,\tfrac19\Bigr], \quad \Bigl[\tfrac29,\tfrac39\Bigr], \quad
   \Bigl[\tfrac69,\tfrac79\Bigr], \quad \Bigl[\tfrac89,1\Bigr]. \]
Continuing in this way, we obtain a sequence of compact sets $E_n$ such that
\begin{enumerate}[label=(\alph*), leftmargin=2.2em, itemsep=1pt, topsep=3pt]
  \item $E_1 \supset E_2 \supset E_3 \supset \cdots$;
  \item $E_n$ is the union of $2^n$ intervals, each of length $3^{-n}$.
\end{enumerate}
The set
\[ P = \bigcap_{n=1}^{\infty} E_n \]
is called the \emph{Cantor set}. $P$ is clearly compact, and Theorem 2.36 shows
that $P$ is not empty.

\begin{center}
\begin{tikzpicture}[x=1mm, y=1mm]
  \def\W{82}
  % E0
  \draw[seg] (0,0) -- (\W,0);
  \node[lbl, anchor=east] at (-3,0) {$E_0$};
  \node[mlbl, anchor=north] at (0,-2) {$0$};
  \node[mlbl, anchor=north] at (\W,-2) {$1$};
  % E1
  \foreach \a/\b in {0/1, 2/3}{
    \draw[seg] ({\a*\W/3},-7) -- ({\b*\W/3},-7);}
  \node[lbl, anchor=east] at (-3,-7) {$E_1$};
  % E2
  \foreach \a/\b in {0/1, 2/3, 6/7, 8/9}{
    \draw[seg] ({\a*\W/9},-14) -- ({\b*\W/9},-14);}
  \node[lbl, anchor=east] at (-3,-14) {$E_2$};
  % E3
  \foreach \a in {0,2,6,8,18,20,24,26}{
    \draw[seg] ({\a*\W/27},-21) -- ({(\a+1)*\W/27},-21);}
  \node[lbl, anchor=east] at (-3,-21) {$E_3$};
  % E4
  \foreach \a in {0,2,6,8,18,20,24,26,54,56,60,62,72,74,78,80}{
    \draw[draw=RuInk, line width=1.2pt]
      ({\a*\W/81},-27) -- ({(\a+1)*\W/81},-27);}
  \node[lbl, anchor=east] at (-3,-27) {$E_4$};
  \node[lbl, anchor=east] at (-3,-33) {$\vdots$};
  \node[lbl, anchor=west] at (\W+4,-14) {$2^n$ pieces,};
  \node[lbl, anchor=west] at (\W+4,-18) {each $3^{-n}$ long};
\end{tikzpicture}
\end{center}
\figcap{Removing middle thirds. Total length removed is
$\sum 2^{n-1}3^{-n} = 1$, so ``nothing'' is left --- yet what remains is
uncountable.}

No segment of the form
\begin{equation}
  \Bigl(\frac{3k+1}{3^m},\ \frac{3k+2}{3^m}\Bigr) \tag{2.24}
\end{equation}
where $k$ and $m$ are positive integers, has a point in common with $P$. Since
every segment $(\alpha,\beta)$ contains a segment of the form (2.24), if
\[ 3^{-m} < \frac{\beta-\alpha}{6}, \]
$P$ contains no segment.

To show that $P$ is perfect, it is enough to show that $P$ contains no isolated
point. Let $x \in P$, and let $S$ be any segment containing $x$. Let $I_n$ be
that interval of $E_n$ which contains $x$. Choose $n$ large enough, so that
$I_n \subset S$. Let $x_n$ be an endpoint of $I_n$, such that $x_n \ne x$.

It follows from the construction of $P$ that $x_n \in P$. Hence $x$ is a limit
point of $P$, and $P$ is perfect.\kn{The endpoints are the key: they are never
removed at any later stage, because removal always takes an \emph{open} middle
third. So every $x \in P$ has other points of $P$ arbitrarily close --- the
endpoints of the shrinking interval containing it.}

One of the most interesting properties of the Cantor set is that it provides us
with an example of an uncountable set of measure zero (the concept of measure
will be discussed in Chap.\ 11).

\begin{deeper}[Everything the Cantor set breaks]
It is worth listing what $P$ demonstrates, because each item contradicts a
plausible belief.

\smallskip
\textbf{Uncountable but containing no interval.} By 2.43 it is uncountable,
being perfect and nonempty; yet it contains no segment. So ``uncountable''
carries no implication of thickness.

\smallskip
\textbf{Total length zero.} The removed intervals have total length
$\frac13 + \frac29 + \frac{4}{27} + \cdots = 1$. What remains has measure zero
(Chapter 11) --- an uncountable set of measure zero.

\smallskip
\textbf{Every point is a limit of \emph{other} points, yet no point has
neighbours in bulk.} $P$ is perfect, so no isolated points; but it is also
nowhere dense, so its closure has empty interior. Both at once.

\smallskip
\textbf{It is not a curiosity.} $P$ is homeomorphic to $\{0,1\}^\N$, which is
why 2.14 and 2.43 measure the same thing --- a point of $P$ is exactly a choice
of ``left third or right third'' at every stage, i.e.\ a $0$--$1$ sequence.
That correspondence also makes the uncountability immediate without 2.43.

\medskip
\begin{center}
\begin{tikzpicture}[x=1mm, y=1mm]
  % level 0
  \draw[seg] (0,16) -- (36,16);
  % level 1
  \draw[seg] (0,8) -- (12,8);   \draw[seg] (24,8) -- (36,8);
  \draw[thin] (16,15) -- (7,9.5);   \node[mlbl, anchor=east] at (11,13) {$0$};
  \draw[thin] (20,15) -- (29,9.5);  \node[mlbl, anchor=west] at (25.5,13) {$1$};
  % level 2
  \draw[seg] (0,0) -- (4,0);   \draw[seg] (8,0) -- (12,0);
  \draw[seg] (24,0) -- (28,0); \draw[seg] (32,0) -- (36,0);
  \draw[thin] (4.5,7) -- (2,1.5);   \node[mlbl, anchor=east] at (3,5) {$0$};
  \draw[thin] (7.5,7) -- (10,1.5);  \node[mlbl, anchor=west] at (9,5) {$1$};
  \draw[thin] (28.5,7) -- (26,1.5); \node[mlbl, anchor=east] at (27,5) {$0$};
  \draw[thin] (31.5,7) -- (34,1.5); \node[mlbl, anchor=west] at (33,5) {$1$};
  \node[mlbl] at (18,-4) {$\vdots$};
  \node[lbl, anchor=west] at (44,10) {a point of $P$ $=$ an infinite};
  \node[lbl, anchor=west] at (44,6.5) {path $=$ its address, e.g.\ $0,1,1,0,\dots$};
  \node[lbl, anchor=west] at (44,3) {--- exactly a $0$--$1$ sequence (2.14)};
\end{tikzpicture}
\end{center}
\figcap{Left third or right third, forever: the surviving intervals form a
binary tree, and a point of $P$ is an infinite path through it.}
\end{deeper}

\rsection{Connected Sets}

\begin{signpost}[Two definitions and one theorem]
The shortest section in the chapter, and it exists for the intermediate value
theorem (4.23), which will follow in one line from Theorem 4.22: continuous
images of connected sets are connected.

\smallskip
The only subtlety is that ``connected'' is defined via \emph{separated} sets,
not disjoint ones. Remark 2.46 is there to stop you conflating the two, and
2.47 shows that on the line the whole notion collapses to something familiar:
connected means interval.
\end{signpost}

\ritem{2.45}{Definition}
Two subsets $A$ and $B$ of a metric space $X$ are said to be \emph{separated}
if both $A \cap \overline{B}$ and $\overline{A} \cap B$ are empty, i.e., if no
point of $A$ lies in the closure of $B$ and no point of $B$ lies in the closure
of $A$.

A set $E \subset X$ is said to be \emph{connected} if $E$ is not a union of two
nonempty separated sets.

\ritem{2.46}{Remark}
Separated sets are of course disjoint, but disjoint sets need not be separated.
For example, the interval $[0,1]$ and the segment $(1,2)$ are not separated,
since $1$ is a limit point of $(1,2)$. However, the segments $(0,1)$ and
$(1,2)$ are separated.

\begin{center}
\begin{tikzpicture}[x=1mm, y=1mm]
  % disjoint but not separated
  \draw[axis] (-3,8) -- (76,8);
  \draw[seg] (6,8) -- (34,8);
  \node[ptfill, minimum size=2.6pt] at (6,8) {};
  \node[ptfill, minimum size=2.6pt] at (34,8) {};
  \draw[seg] (34,8) -- (62,8);
  \node[ptopen] at (34,8) {};
  \node[ptopen] at (62,8) {};
  \node[mlbl, anchor=north] at (6,6.4) {$0$};
  \node[mlbl, anchor=north] at (34,6.4) {$1$};
  \node[mlbl, anchor=north] at (62,6.4) {$2$};
  \node[lbl, anchor=west] at (80,8) {disjoint, \emph{not} separated};
  % separated
  \draw[axis] (-3,-8) -- (76,-8);
  \draw[seg] (6,-8) -- (34,-8);
  \node[ptopen] at (6,-8) {}; \node[ptopen] at (34,-8) {};
  \draw[seg] (34,-8) -- (62,-8);
  \node[ptopen] at (34,-8) {}; \node[ptopen] at (62,-8) {};
  \node[mlbl, anchor=north] at (34,-9.6) {$1$};
  \node[lbl, anchor=west] at (80,-8) {separated};
\end{tikzpicture}
\end{center}
\figcap{The difference is whether the shared boundary point belongs to one of
the sets. Closures must miss each other, not merely the sets themselves.}

\begin{pitfall}[Separated is stronger than disjoint, and the gap matters]
If ``connected'' were defined by disjointness, \emph{every} set with more than
one point would be disconnected --- split $[0,1]$ into $[0,\frac12]$ and
$(\frac12,1]$ and you have two disjoint nonempty pieces. The definition would
be useless.

Requiring the \emph{closures} to miss is what rules this out: $\frac12$ lies in
the closure of the second piece, so the split is not a separation. In effect,
separated sets are ones that do not touch even in the limit.
\end{pitfall}

The connected subsets of the line have a particularly simple structure:

\ritem{2.47}{Theorem}
\emph{A subset $E$ of the real line $\R^1$ is connected if and only if it has
the following property: If $x \in E$, $y \in E$, and $x < z < y$, then
$z \in E$.}

\begin{proof}
If there exist $x \in E$, $y \in E$, and some $z \in (x,y)$ such that
$z \notin E$, then $E = A_z \cup B_z$ where
\[ A_z = E \cap (-\infty,z), \qquad B_z = E \cap (z,\infty). \]
Since $x \in A_z$ and $y \in B_z$, $A$ and $B$ are nonempty. Since
$A_z \subset (-\infty,z)$ and $B_z \subset (z,\infty)$, they are separated.
Hence $E$ is not connected.

To prove the converse, suppose $E$ is not connected. Then there are nonempty
separated sets $A$ and $B$ such that $A \cup B = E$. Pick $x \in A$,
$y \in B$, and assume (without loss of generality) that $x < y$. Define
\[ z = \sup\bigl(A \cap [x,y]\bigr). \]
By Theorem 2.28, $z \in \overline{A}$; hence $z \notin B$. In particular,
$x \le z < y$.\kn{Chapter 1 again. The supremum exists by 1.19, lands in
$\overline A$ by 2.28, and therefore avoids $B$ because $A$ and $B$ are
separated --- that is the one place separation is spent. Every step of this
converse is an application of something already proved.}

If $z \notin A$, it follows that $x < z < y$ and $z \notin E$.

If $z \in A$, then $z \notin \overline{B}$, hence there exists $z_1$ such that
$z < z_1 < y$ and $z_1 \notin B$. Then $x < z_1 < y$ and $z_1 \notin E$.
\end{proof}

\begin{center}
\begin{tikzpicture}[x=1mm, y=1mm]
  \draw[axis] (-2,0) -- (98,0);
  \node[ptfill, minimum size=2.6pt] at (8,0) {};
  \node[mlbl, anchor=north] at (8,-2) {$x\in A$};
  \node[ptfill, minimum size=2.6pt] at (90,0) {};
  \node[mlbl, anchor=north] at (90,-2) {$y\in B$};
  \draw[seg] (8,0) -- (46,0);
  \node[lbl, anchor=north] at (27,-3.5) {$A\cap[x,y]$};
  \draw[guide] (46,-7) -- (46,12);
  \node[lbl, anchor=south] at (46,12) {$z=\sup\bigl(A\cap[x,y]\bigr)$};
  \node[ptopen] at (54,0) {};
  \node[mlbl, anchor=north] at (54,-2) {$z_1$};
  \node[lbl, anchor=west] at (58,10.5) {$z\in\overline{A}$, so $z\notin B$:};
  \node[lbl, anchor=west] at (58,7) {either $z\notin E$ already, or room};
  \node[lbl, anchor=west] at (58,3.5) {just past $z$ misses $B$ too};
\end{tikzpicture}
\end{center}
\figcap{The same seam picture as Definition 1.8, one chapter later: the
supremum lands in $\overline{A}$ and hence off $B$, and a point missing from
$E$ is found at or just past the seam.}

\begin{unpack}[Connected on the line means interval]
The property in 2.47 --- if two points are in $E$, everything between them is
too --- is exactly the definition of an interval, in the general sense that
includes $\emptyset$, single points, half-lines and $\R$ itself. So the theorem
reads: \emph{the connected subsets of $\R$ are precisely the intervals.}

\smallskip
That is worth holding because it makes the intermediate value theorem
(4.23) transparent. If $f$ is continuous
on $[a,b]$ then $f([a,b])$ is connected (4.22), hence an interval by 2.47,
hence contains every value between $f(a)$ and $f(b)$ --- the intermediate value
theorem. Chapter 1's Theorem 1.21 on $n$th roots was a hand-rolled special case
of exactly this, as the note there observed.

\smallskip
Note also that connectedness, unlike openness, is intrinsic: like compactness
it survives change of ambient space, since 2.45 refers only to closures within
$E$'s own metric.
\end{unpack}

\begin{recap}[Chapter 2 in one paragraph]
Countability sorts infinite sets into two sizes, and $\Q$ and $\R$ land on
opposite sides (2.13, 2.43). Metric spaces distill $\R^k$ down to its distance
function, and open, closed and limit point are the vocabulary that distance
alone supports (2.15--2.30). Compactness --- every open cover admits a finite
subcover --- is the property that turns local facts into global ones, and in
$\R^k$ it is exactly closed-and-bounded (2.41), a fact that rests on
completeness through the nested-interval theorems (2.38--2.40) and on nothing
else. Perfect sets are uncountable (2.43), and the Cantor set shows how thin
one can be (2.44). Connected subsets of the line are precisely the intervals
(2.47). Chapter 4 will spend all of it: continuity gets defined by inverse
images of open sets, compactness becomes boundedness and attainment of extrema,
and connectedness becomes the intermediate value theorem.
\end{recap}
